% GENERATED FILE. Edit canonical lessons, then run npm run generate:books.
% canonical-source-hash: fnv1a64:8838a8b03b22651e
% canonical-lessons: SA-C181-dantapida, SA-C181-pratisyayah, SA-C181-svasah, SA-C181-nadi, SA-C181-svasthyam, SA-R181-first-pass-words-from-chapters-173-176, SA-R181-second-pass-words-from-chapters-177-181

\chapter{Toothache, Head cold, Breath, Pulse, Fitness}
\label{ch:sa-toothache-head-cold-breath-pulse-fitness}
\hlchaptermodality{181}

\begin{chapteropening}
\textbf{By the end of this chapter:} \emph{I can say a toothache, a head cold, breath, the pulse and fitness.}

Complete the last lesson of chapter 181: I can say a toothache, a head cold, breath, the pulse and fitness.
\end{chapteropening}

\section[\texorpdfstring{\sk{दन्तपीडा} (dantapīḍā)}{दन्तपीडा (dantapīḍā)}]{\sk{दन्तपीडा} (dantapīḍā) — a toothache}

\label{lesson:SA-C181-dantapida}



\begin{quote}
Before the new one: say the Sanskrit for a scholar, then the Sanskrit for a learned person.
\end{quote}

\subsection*{You'll want to know: \sk{दन्तपीडा}}
\textbf{\sk{दन्तपीडा}} — \emph{dantapīḍā} — \textquotedblleft{}a toothache\textquotedblright{}.

\begin{grammarlens}[title={What you've built}]
One more word for this chapter.
\end{grammarlens}

\subsection*{Your turn}
\begin{itemize}
\raggedright
  \item \emph{Say it:} \emph{dantapīḍā}
  \item \emph{Say it:} \emph{dantapīḍā}, once more
  \item \emph{Say it:} say \emph{dantapīḍā}
  \item \emph{Recall:} say the Sanskrit for a scholar, then the Sanskrit for a learned person, then say \emph{dantapīḍā} again
\end{itemize}

\begin{tcolorbox}[breakable,colback=teal!4,colframe=teal!35!black,title={Before you move on}]
What does \sk{दन्तपीडा} mean? (\textquotedblleft{}A toothache\textquotedblright{}.) Say it once more.
\end{tcolorbox}
\section[\texorpdfstring{\sk{प्रतिश्यायः} (pratiśyāyaḥ)}{प्रतिश्यायः (pratiśyāyaḥ)}]{\sk{प्रतिश्यायः} (pratiśyāyaḥ) — a head cold}

\label{lesson:SA-C181-pratisyayah}



\begin{quote}
Before the new one: say the Sanskrit for a learned person, then the Sanskrit for a toothache.
\end{quote}

\subsection*{You'll want to know: \sk{प्रतिश्यायः}}
\textbf{\sk{प्रतिश्यायः}} — \emph{pratiśyāyaḥ} — \textquotedblleft{}a head cold\textquotedblright{}.

\begin{grammarlens}[title={What you've built}]
One more word for this chapter.
\end{grammarlens}

\subsection*{Your turn}
\begin{itemize}
\raggedright
  \item \emph{Say it:} \emph{pratiśyāyaḥ}
  \item \emph{Say it:} \emph{pratiśyāyaḥ}, once more
  \item \emph{Say it:} say \emph{pratiśyāyaḥ}
  \item \emph{Recall:} say the Sanskrit for a learned person, then the Sanskrit for a toothache, then say \emph{pratiśyāyaḥ} again
\end{itemize}

\begin{tcolorbox}[breakable,colback=teal!4,colframe=teal!35!black,title={Before you move on}]
What does \sk{प्रतिश्यायः} mean? (\textquotedblleft{}A head cold\textquotedblright{}.) Say it once more.
\end{tcolorbox}
\section[\texorpdfstring{\sk{श्वासः} (śvāsaḥ)}{श्वासः (śvāsaḥ)}]{\sk{श्वासः} (śvāsaḥ) — breath}

\label{lesson:SA-C181-svasah}



\begin{quote}
Before the new one: say the Sanskrit for a toothache, then the Sanskrit for a head cold.
\end{quote}

\subsection*{You'll want to know: \sk{श्वासः}}
\textbf{\sk{श्वासः}} — \emph{śvāsaḥ} — \textquotedblleft{}breath\textquotedblright{}.

\begin{grammarlens}[title={What you've built}]
One more word for this chapter.
\end{grammarlens}

\subsection*{Your turn}
\begin{itemize}
\raggedright
  \item \emph{Say it:} \emph{śvāsaḥ}
  \item \emph{Say it:} \emph{śvāsaḥ}, once more
  \item \emph{Say it:} say \emph{śvāsaḥ}
  \item \emph{Recall:} say the Sanskrit for a toothache, then the Sanskrit for a head cold, then say \emph{śvāsaḥ} again
\end{itemize}

\begin{tcolorbox}[breakable,colback=teal!4,colframe=teal!35!black,title={Before you move on}]
What does \sk{श्वासः} mean? (\textquotedblleft{}Breath\textquotedblright{}.) Say it once more.
\end{tcolorbox}
\section[\texorpdfstring{\sk{नाडी} (nāḍī)}{नाडी (nāḍī)}]{\sk{नाडी} (nāḍī) — the pulse}

\label{lesson:SA-C181-nadi}



\begin{quote}
Before the new one: say the Sanskrit for a head cold, then the Sanskrit for breath.
\end{quote}

\subsection*{You'll want to know: \sk{नाडी}}
\textbf{\sk{नाडी}} — \emph{nāḍī} — \textquotedblleft{}the pulse\textquotedblright{}.

\begin{grammarlens}[title={What you've built}]
One more word for this chapter.
\end{grammarlens}

\subsection*{Your turn}
\begin{itemize}
\raggedright
  \item \emph{Say it:} \emph{nāḍī}
  \item \emph{Say it:} \emph{nāḍī}, once more
  \item \emph{Say it:} say \emph{nāḍī}
  \item \emph{Recall:} say the Sanskrit for a head cold, then the Sanskrit for breath, then say \emph{nāḍī} again
\end{itemize}

\begin{tcolorbox}[breakable,colback=teal!4,colframe=teal!35!black,title={Before you move on}]
What does \sk{नाडी} mean? (\textquotedblleft{}The pulse\textquotedblright{}.) Say it once more.
\end{tcolorbox}
\section[\texorpdfstring{\sk{स्वास्थ्यम्} (svāsthyam)}{स्वास्थ्यम् (svāsthyam)}]{\sk{स्वास्थ्यम्} (svāsthyam) — fitness, health}

\label{lesson:SA-C181-svasthyam}



\begin{quote}
Before the new one: say the Sanskrit for breath, then the Sanskrit for the pulse.
\end{quote}

\subsection*{You'll want to know: \sk{स्वास्थ्यम्}}
\textbf{\sk{स्वास्थ्यम्}} — \emph{svāsthyam} — \textquotedblleft{}fitness, health\textquotedblright{}.

\begin{grammarlens}[title={What you've built}]
One more word for this chapter.
\end{grammarlens}

\subsection*{Your turn}
\begin{itemize}
\raggedright
  \item \emph{Say it:} \emph{svāsthyam}
  \item \emph{Say it:} \emph{svāsthyam}, once more
  \item \emph{Say it:} say \emph{svāsthyam}
  \item \emph{Recall:} say the Sanskrit for breath, then the Sanskrit for the pulse, then say \emph{svāsthyam} again
\end{itemize}

\begin{tcolorbox}[breakable,colback=teal!4,colframe=teal!35!black,title={Before you move on}]
What does \sk{स्वास्थ्यम्} mean? (\textquotedblleft{}Fitness, health\textquotedblright{}.) Say it once more.
\end{tcolorbox}
\section[(dialogue)]{Words from chapters 173-176 — a first pass}

\label{lesson:SA-R181-first-pass-words-from-chapters-173-176}



\begin{quote}
Say the words for the pulse and fitness.
\end{quote}

\subsection*{Your turn}
Say these aloud:

\begin{itemize}
\raggedright
  \item throws, lets go, gets, serves, enjoys
  \item steals, eats up, strikes, keeps, hurts
  \item takes away, supports, carries, endures, tries
  \item changes, reports, invites, cleans, fills
\end{itemize}

\begin{tcolorbox}[breakable,colback=teal!4,colframe=teal!35!black,title={Before you move on}]
Throws? (\textbf{\sk{क्षिपति}}, \emph{kṣipati}.) Lets go? (\textbf{\sk{मुञ्चति}}, \emph{muñcati}.) Gets? (\textbf{\sk{लभते}}, \emph{labhate}.) Serves? (\textbf{\sk{सेवते}}, \emph{sevate}.) Enjoys? (\textbf{\sk{रमते}}, \emph{ramate}.) Steals? (\textbf{\sk{चोरयति}}, \emph{corayati}.) Eats up? (\textbf{\sk{भक्षयति}}, \emph{bhakṣayati}.) Strikes? (\textbf{\sk{ताडयति}}, \emph{tāḍayati}.) Keeps? (\textbf{\sk{पालयति}}, \emph{pālayati}.) Hurts? (\textbf{\sk{पीडयति}}, \emph{pīḍayati}.) Takes away? (\textbf{\sk{हरति}}, \emph{harati}.) Supports? (\textbf{\sk{भरति}}, \emph{bharati}.) Carries? (\textbf{\sk{वहति}}, \emph{vahati}.) Endures? (\textbf{\sk{सहते}}, \emph{sahate}.) Tries? (\textbf{\sk{यतते}}, \emph{yatate}.) Changes? (\textbf{\sk{परिवर्तयति}}, \emph{parivartayati}.) Reports? (\textbf{\sk{निवेदयति}}, \emph{nivedayati}.) Invites? (\textbf{\sk{आमन्त्रयति}}, \emph{āmantrayati}.) Cleans? (\textbf{\sk{शोधयति}}, \emph{śodhayati}.) Fills? (\textbf{\sk{पूरयति}}, \emph{pūrayati}.)
\end{tcolorbox}
\section[(dialogue)]{Words from chapters 177-181 — a second pass}

\label{lesson:SA-R181-second-pass-words-from-chapters-177-181}



\begin{quote}
Say the words for conceals and distributes.
\end{quote}

\subsection*{Your turn}
Say these aloud:

\begin{itemize}
\raggedright
  \item conceals, distributes, collects, trusts, is ashamed
  \item a mother's brother's wife, a father-in-law, a mother-in-law, a son-in-law, a daughter-in-law
  \item a soldier, a guard, a judge, a lawyer, a writer
  \item a hunter, a thief, a beggar, a scholar, a learned person
  \item a toothache, a head cold, breath, the pulse, fitness
\end{itemize}

\begin{tcolorbox}[breakable,colback=teal!4,colframe=teal!35!black,title={Before you move on}]
Conceals? (\textbf{\sk{प्रच्छादयति}}, \emph{pracchādayati}.) Distributes? (\textbf{\sk{वितरति}}, \emph{vitarati}.) Collects? (\textbf{\sk{संगृह्णाति}}, \emph{saṅgṛhṇāti}.) Trusts? (\textbf{\sk{विश्वसिति}}, \emph{viśvasiti}.) Is ashamed? (\textbf{\sk{लज्जते}}, \emph{lajjate}.) A mother's brother's wife? (\textbf{\sk{मातुलानी}}, \emph{mātulānī}.) A father-in-law? (\textbf{\sk{श्वशुरः}}, \emph{śvaśuraḥ}.) A mother-in-law? (\textbf{\sk{श्वश्रूः}}, \emph{śvaśrūḥ}.) A son-in-law? (\textbf{\sk{जामाता}}, \emph{jāmātā}.) A daughter-in-law? (\textbf{\sk{स्नुषा}}, \emph{snuṣā}.) A soldier? (\textbf{\sk{सैनिकः}}, \emph{sainikaḥ}.) A guard? (\textbf{\sk{रक्षकः}}, \emph{rakṣakaḥ}.) A judge? (\textbf{\sk{न्यायाधीशः}}, \emph{nyāyādhīśaḥ}.) A lawyer? (\textbf{\sk{अधिवक्ता}}, \emph{adhivaktā}.) A writer? (\textbf{\sk{लेखकः}}, \emph{lekhakaḥ}.) A hunter? (\textbf{\sk{व्याधः}}, \emph{vyādhaḥ}.) A thief? (\textbf{\sk{चौरः}}, \emph{cauraḥ}.) A beggar? (\textbf{\sk{याचकः}}, \emph{yācakaḥ}.) A scholar? (\textbf{\sk{पण्डितः}}, \emph{paṇḍitaḥ}.) A learned person? (\textbf{\sk{विद्वान्}}, \emph{vidvān}.) A toothache? (\textbf{\sk{दन्तपीडा}}, \emph{dantapīḍā}.) A head cold? (\textbf{\sk{प्रतिश्यायः}}, \emph{pratiśyāyaḥ}.) Breath? (\textbf{\sk{श्वासः}}, \emph{śvāsaḥ}.) The pulse? (\textbf{\sk{नाडी}}, \emph{nāḍī}.) Fitness? (\textbf{\sk{स्वास्थ्यम्}}, \emph{svāsthyam}.)
\end{tcolorbox}
